% generated by src/analysis/report.py --- do not edit by hand
\begin{table}[htbp]\centering
  \caption{Factor ablation for detection (\texttt{pDetA}): McFadden pseudo-$R^2$ and the out-of-sample MAE gain over a mean baseline under leave-species-out (LSO) and leave-location-out (LLO). Dropping any single distance barely changes the fit or the (null) out-of-sample gain --- no factor carries signal; only adding the animal-size covariate moves the score.}
  \label{tab:ablation}
  \begin{tabular}{lccc}
    \toprule
    Design & pseudo-$R^2$ & LSO $\Delta$~($p$) & LLO $\Delta$~($p$) \\
    \midrule
    full (4 distances) & 0.055 & $+0.004$~($p{=}0.24$) & $+0.004$~($p{=}0.23$) \\
    drop taxonomic & 0.039 & $+0.000$~($p{=}0.47$) & $+0.000$~($p{=}0.49$) \\
    drop temporal & 0.055 & $+0.004$~($p{=}0.24$) & $+0.004$~($p{=}0.23$) \\
    drop visual & 0.048 & $+0.004$~($p{=}0.22$) & $+0.004$~($p{=}0.28$) \\
    drop environment & 0.055 & $+0.006$~($p{=}0.15$) & $+0.005$~($p{=}0.17$) \\
    $+$ animal size & 0.119 & $+0.018$~($p{=}0.02$) & $+0.020$~($p{=}0.00$) \\
    $-$ support covariate & 0.035 & $-0.002$~($p{=}0.72$) & $-0.003$~($p{=}0.73$) \\
    \bottomrule
  \end{tabular}
\end{table}
